\setauthor{\secondauthor}
\section{Personalvermittlung in der Gastronomie}

Die Gastronomie zeichnet sich durch eine hohe Dynamik des Arbeitsmarktes aus.
Der Bedarf an Arbeitskräften entsteht oft kurzfristig, etwa durch saisonale Schwankungen,
Veranstaltungen, Ausfälle, bedingt durch Krankheit oder wechselnde Auslastungen.
Gleichzeitig kämpft die Branche seit Jahren mit strukturellen Schwierigkeiten,
wie instabilen Arbeitsbedingungen, hoher Fluktuation und 
Rekrutierungsschwierigkeiten. Giousmpasoglou beschreibt in seiner aktuellen 
Übersicht zu den Arbeitsbedingungen in der Hospitality-Branche unter anderem 
Einkommensinstabilität, Unterbeschäftigung, hohen Arbeitsdruck sowie gravierende 
Schwierigkeiten bei Recruiting und Mitarbeiterbindung \cite{Giousmpasoglou2024}.

Für die Personalvermittlung heißt das, dass die traditionellen und informellen 
Wege der Rekrutierung zwar weiterhin in vielen Fällen beschritten werden, aber 
nur noch eingeschränkt den Anforderungen eines zunehmend flexibleren 
Arbeitsmarktes gerecht werden. Gerade kurzfristige Einsätze verlangen eine 
schnelle, transparente und beiden Seiten leicht zugängliche Vermittlungsform. 
Eine zunehmende Orientierung auf die sogenannte Gig Economy zeigt sich 
insbesondere in der Hospitality-Branche. El Hajal und Rowson weisen darauf hin, 
dass flexible, kurzfristige und plattformvermittelte Arbeitsformen im 
Hospitality-Kontext an Bedeutung gewinnen und die klassische Struktur der 
Beschäftigungsverhältnisse ergänzen oder teilweise ablösen
\cite{ElHajal2021}.

Damit wird deutlich: Digitale Vermittlungsplattformen sind nicht nur 
technische Spielerei, sondern eine echte Branchenanliegen, die helfen können, 
Informationen über Verfügbarkeit, Einsatzmöglichkeiten und Anforderungen zu 
bündeln und den Vermittlungsprozess für beide Seiten einfacher zu handhaben. 
Dabei ist für die vorliegende Arbeit besonders wichtig, dass mit einer digitalen 
Plattform allein noch keine erfolgreiche Vermittlung sichergestellt ist. Damit 
ein solches System im Markt funktionieren kann, muss es erst einmal potenzielle 
Nutzer:innen davon überzeugen, die Anwendung wahrzunehmen, ihr zu vertrauen und 
sie tatsächlich zu nutzen.

\section{Marketing-Chatbots}

Bei Marketing-Chatbots handelt es sich um dialogbasierte Systeme, die aktive 
Eingriffe in Kommunikations- und Entscheidungsprozesse durchführen können. 
Im Unterschied zu herkömmlichen Webseiten stellen sie 
Inhalte nicht nur statisch und linear, sondern interaktiv im Gespräch dar. Dies 
ist für das Marketing von Bedeutung, weil Inhalte an die Situation angepasst werden
können und Nutzer:innen schrittweise durch den Informationsprozess geführt werden
können.

Die zunehmende Bedeutung von Conversational Agents im wirtschaftlichen Kontext
wird durch mehrere Übersichtsarbeiten bestätigt. In ihrer systematischen
Literaturübersicht kommen Bavaresco et al. zu dem Ergebnis, dass Unternehmen 
Chatbots vermehrt einsetzen, um die Servicequalität zu erhöhen, Kosten zu sparen 
und die Interaktion mit Kund:innen zu erleichtern \cite{Bavaresco2020}.
Ling et al. weisen in ihrer systematischen
Review zur Nutzung und Akzeptanz von Conversational Agents darauf hin, dass
der Erfolg der Implementierung solcher Systeme stark davon abhängen wie die 
Nutzungsvorteile wahrgenommen werden \cite{Ling2021}.
Für Marketing-Anwendungen heißt das, dass Chatbots
vor allem dann von Bedeutung sind, wenn sie von Nutzer:innen als hilfreich und
verständlich erlebt werden.

Für die konkrete Nutzererfahrung mit Chatbots sind die Ergebnisse von Følstad und 
Brandtzaeg wichtig. Die Autoren zeigen, dass positive Nutzererfahrungen vor allem 
mit gewissen pragmatischen Faktoren wie effektiver Hilfeleistung und 
Problemlösungen verbunden sind. Negative Erfahrungen ergeben sich dagegen häufig 
durch Fehlinterpretationen, unpassende Antworten oder irritierendes Verhalten \cite{Folstad2020}.
Für die Entwicklung eines Marketing-Chatbots folgt daraus, dass der Neuheitswert 
eines Chatbots nicht genügt. Der Dialog muss so gestaltet sein, dass Informationen 
zielgerichtet und ohne unnötige Frustration und kognitive Belastung vermittelt 
werden.

Zusätzlich sollten die oben genannten Ergebnisse auch
durch hemmende Faktoren relativiert werden. Jan et al. nennen für
shoppingbezogene Conversational Agents als Gründe für die Nutzung u.a.
Nützlichkeit, Informationsgehalt und Interaktivität, dagegen aber auch
funktionale Risiken und perceived intrusiveness als Gründe gegen die
Nutzung \cite{Jan2023}. Diese Ergebnisse sind für unser Projekt von Relevanz, da der
Chatbot ja nicht nur Informationen bereitstellen, sondern auch Vertrauen
schaffen und die Nutzer:innen zu einem App-Download führen soll.

\section{Personalisierung und Tonalität von Chatbots}

Ein weiterer entscheidender Faktor für die Gestaltung von Marketing-Chatbots betrifft
die Frage, wie stark die Inhalte personalisiert und wie sie sprachlich
ausgestaltet werden sollen. Personalisierung kann in einem
unternehmensspezifischen Kontext beispielsweise durch die Berücksichtigung von
Nutzerrollen, das Anpassen des Beispiels an Nutzer:innen, die Verwendung des Namens oder auch
einen variierenden Gesprächsverlauf stattfinden. Grundannahme ist, dass
Inhalte wirksamer sind, wenn sie von der Person, an die sie gerichtet sind,
als relevant für sie wahrgenommen werden.

Diese Annahme hat aktuellen Befunden zufolge ihre Richtigkeit. In einer
Meta-Analyse auf Basis von 53 experimentellen Studien zeigen Yeo, Chu und
Li, dass personalisierte Werbung im Mittel wirksamer ist als generische
Werbung und insbesondere positive Effekte auf Einstellungen und
Verhaltensintentionen hat. Der offenbar wirksamste Mechanismus ist dabei die
erhöhte wahrgenommene Relevanz der Botschaft 
\cite{Yeo2025meta}. Die Untersuchung bezieht sich
zwar auf personalisierte Werbung im Allgemeinen, dennoch ist der Befund für
die Gestaltung digitaler Informationssysteme direkt zutreffend: Wenn
Informationen enger an die Situation einer Zielgruppe angepasst sind, ist die
Chance größer, dass sie von dieser auch als bedeutsam und überzeugend
erlebt werden.

Die Rolle der Personalisierung ist auch für Conversational Agents bereits 
systematisch untersucht worden. In einer systematischen Review zeigen 
Kocaballi et al. , dass Personalisierung in Chatbots mit höherer 
Nutzerzufriedenheit, stärkerem Engagement, besserer Dialogqualität usw. in 
Verbindung steht. Gleichzeitig stellen die Autoren fest, dass Personalisierung in 
vielen bestehenden Systemen zwar genutzt, aber nur wenig systematisch evaluiert 
wird \cite{Kocaballi2019}. Daraus ist für diese Arbeit abzuleiten, dass Personalisierung nicht nur 
verbaut, sondern auch gezielt als eigenen Untersuchungsfaktor betrachtet werden 
sollte.


Neben der Personalisierung des Inhalts ist auch die Tonalität des Chatbots von 
Bedeutung. Sie bestimmt, ob das System als formell oder seriös, nahbar oder 
freundlich wahrgenommen wird. Gerade im Marketing spielt dieser Aspekt eine Rolle, 
da sprachliche Gestaltung nicht nur Information bietet, sondern auch zur 
Markenwirkung und Vertrauensbildung beiträgt. Følstad und Brandtzaeg zeigen, dass 
neben den pragmatischen auch hedonische Aspekte das Nutzererlebnis mit Chatbots 
bestimmen, darunter Unterhaltung, Natürlichkeit der Formulierungen und das 
Ausbleiben irritierender oder unangebrachter Antworten. Für das Projekt wird 
damit deutlich, dass Tonalität nicht als eine rein stilistische Frage sondern als 
Gestaltungsparameter mit Auswirkung auf Akzeptanz und Zufriedenheit zu verstehen 
ist \cite{Folstad2020}.

Die theoretische Herleitung hat für die Diplomarbeit die Konsequenz, dass ein 
Experiment zur Untersuchung des Einflusses von Personalisierung und Tonalität auf 
Nutzerverhalten und Nutzerwahrnehmung nicht als Verfeinerung, sondern als 
zentrale Variabeln angelegt werden kann.


\section{Tracking und Analyse von Nutzerverhalten}

Um verschiedene Informationsstrategien vergleichen zu können, ist es wichtig, 
das Nutzerverhalten zu erfassen. Im digitalen Marketing wird diese Aufgabe 
normalerweise durch Web- und Marketing-Analytics übernommen. Dabei geht es nicht 
nur um Besucherzahlen, sondern vor allem um das Verhalten entlang eines Funnels 
einer Interaktionskette: welche Schritte werden erreicht, wo gibt es Brüche, 
welche Inhalte werden besonders viel genutzt, und wo gibt es Signale für 
Konversionsbereitschaft?

In ihrer offenen Forschungsarbeit zu Web Analytics und Sales Funnels zeigt 
Golik-Górecka, wie wichtig Web Analytics als Datenbasis für die 
Marketingautomatisierung und auch zur Bewertung von Funnel-Prozessen ist. 
Dabei ist besonders wichtig, dass nicht Kennzahlen isoliert betrachtet werden, 
sondern dass Metriken zu den jeweiligen Funnel-Stufen in Beziehung gesetzt werden \cite{GolikGorecka2023}.
Für diese Diplomarbeit bedeutet das, dass Interaktionsdaten wie Gesprächsdauer, 
Anzahl der Nachrichten, Feedback, Download-Intention nicht als Rohdaten zu 
verstehen sind, sondern als Indikatoren in einer größeren Nutzerführungslogik.

Auch aus der Perspektive der Conversion-Optimierung ist die
Analyse von Touchpoints und Nutzerverhalten entscheidend. Zimmermann
und Auinger verdeutlichen, dass nur eine datenbasierte Analyse der
Kontaktpunkte und Nutzerreaktionen eine zielgerichtete Verbesserung der
Conversion ermöglicht \cite{Zimmermann2023}.
Übertragen auf das vorliegende Projekt heißt das: sowohl der klassische
Funnel als auch der Chatbot sind nicht nur zu implementieren, sondern
anhand spezifischer Kennzahlen zu bewerten. Nur so ist eine
verlässliche Bewertung der beiden Informationsstrategien möglich.

Außerdem ist für die Bewertung dialogischer Systeme wichtig, dass nicht jede 
Interaktionsform automatisch zu besserer Usability führt. In einer offenen 
Usability-Studie zeigen Iftikhar et al. dass ein konversationelles Formular nicht 
in allen Metriken besser abschnitt als ein Single-Page- oder Multi-Page-Ansatz 
\cite{Iftikhar2021}. 
Der Befund ist für diese Arbeit besonders wertvoll, weil er unterstreicht, dass 
ein Chatbot nicht nur wegen seines dialogischen Charakters besser sein kann. 
Empirische Auswertungen tatsächlicher Nutzungsdaten sind erforderlich.

Für das Projekt bedeutet dies ganz konkret:
Sowohl für den Funnel als auch für den Chatbot müssen nachprüfbare
Metriken definiert werden. Dazu gehören die Anzahl der getauschten
Nachrichten, die Dauer einer Sitzung, wie oft welche bestimmten Fragen
gestellt werden, subjektive Zufriedenheitswerte und Indikatoren
für Downloadinteresse. Die Auswertung dieser Daten bildet die
Basis der späteren Variantenbewertung.

\section{Schlussfolgerungen für die Konzeption}

Die Umfeldanalyse hat für die Gestaltung der im Projekt zu entwickelnden Systeme 
mehrere Konsequüsse. Der Überblick über die Personalvermittlung in der 
Gastronomie macht deutlich, daß insbesondere Flexibilität, schnelle 
Informationsbereitstellung und geringe Interaktionsschwelle Forderungen an eine 
digitale Vermittlungslösung sind \cite{Giousmpasoglou2024,ElHajal2021}.
In der Präsentation der App muss daher deutlich 
werden, welchen konkreten Nutzen die Plattform für die jeweilige Zielgruppe 
stiftet.

Zweitens verdeutlichen die zitierten Studien zu Marketing-Chatbots, dass diese 
insbesondere dann sinnvoll eingesetzt werden, wenn sie Informationen gut 
aufbereitet, verständlich und interaktiv vermitteln \cite{Ling2021,Folstad2020,Jan2023}.
Für die Gestaltung 
bedeutet das, dass der Chatbot im Projekt nicht nur Fragen beantworten, sondern 
Nutzer:innen auch strukturiert durch den Informationsprozess leiten muss. Dabei 
ist zu beachten, dass Unklarheit, Missverständnisse oder unnötige Auskünfte die 
Nutzererfahrung deutlich verschlechtern können.

Drittens macht die Literatur zu Personalisierung deutlich, dass 
zielgruppenspezifische und relevante Kommunikation ein wesentlicher Erfolgsfaktor 
sein kann \cite{Yeo2025meta,Kocaballi2019}. Daraus lässt sich ableiten, dass die Unterscheidung zwischen 
Gastronom:innen und Jobsuchenden nicht nur aus inhaltlicher Sicht sinnvoll sein kann, 
sondern auch eine theoretisch begründete Gestaltungsentscheidung darstellt. 
Darüber hinaus ist die Tonalität als eigene Einflussgröße zu werten.

Zudem zeigt die Forschung zu Web Analytics und Conversion-Optimierung,
dass es nicht genügt, einfach einen Funnel oder einen Chatbot zu implementieren. 
Nur durch systematisches Tracking und die spätere Analyse der Nutzungsdaten kann 
beurteilt werden, welche Variante tatsächlich besser ist \cite{GolikGorecka2023,Zimmermann2023,Iftikhar2021}.

Insgesamt stützt die Umfeldanalyse damit die Entscheidung, im weiteren 
Projektverlauf sowohl einen klassischen Multi-Step-Funnel als auch einen 
dialogbasierenden Chatbot zu entwickeln und beide Ansätze datenbasiert zu 
evaluieren.
